% s-071-p071-m001; original PDF 71, 72, 73, 74, 75
\Aphorism{s-071-p071-m001}{64}
Ich werde wahrscheinlich ein schlechter Erzieher werden, im Sinne unserer heutigen erziehungssüchtigen Civilisation, und wenn ich meine Kinder liebe, so sollte ich sie vor mir warnen, was ich hiermit im Voraus gethan haben will. Soviel oder sowenig glaube ich für mich gutsagen zu können: dass jeder Ausbruch von Eigenwillen, Übermuth, freier wilder Natürlichkeit mich ästhetisch zu sehr erfreuen würde, als dass ich, mit ängstlicher Berechnung der Folgen, mich entschlösse den autoritätslüsternen Exercirmeister und pädagogischen Feldwebel zu spielen. Wo bliebe da der Gleichtritt — ich wollte sagen, die „Pflicht“, dieser Grundbegriff aller Erziehung! Ich werde mir versagen müssen, zu den „Verbesserern der Menschheit“ zu gehören; aber ich will mir nicht versagen, dieser Species Mensch mit der Laterne des Cynikers ins Gesicht zu leuchten. Dass sie heute, im Zeitalter der Machtzersplitterung, sich nicht mehr als Priester und Gesetzgeber, sondern im bescheideneren Massstabe als Lehrer und Erzieher unnütz zu machen pflegen, sei ihnen nicht sonderlich hoch angerechnet. Das Arbeitsfeld ist ihnen verengt worden, aber von ihrer Methode haben sie sich nichts abdingen lassen; sie haben sie nur ein wenig modernisirt, verfeinert, sie nehmen nicht mehr das glühende Eisen und die Stachelpeitsche, um die Bestie zu zähmen. Aber die Methode heisst heute noch: ausrotten, abdämmen, beschneiden, unterbinden, einschränken, verbieten — es ist eine grundsätzlich negative, privative, prohibitive Methode des Erziehens, Besserns, Strafens, ein Abschaffen an Stelle des Schaffens, ein Amputiren an Stelle des Heilens, kurz ausgedrückt: wenn dich deines Nächsten Auge ärgert, reisse es aus! Das Sonderbare ist, dass der Erfolg immer für diese Künstler der Zerstörung zu sprechen scheint; wenn die Natur sich willig unterdrücken und vergewaltigen, in ihrer Sprache „verbessern“, lässt, oder wenn sie reich und positiv genug ist, ihren Eingriffen zu trotzen, oder wenn sie sich gar durch excessive Ausbildung einer Anlage hilft, die jenen Quälgeistern entgangen ist —: in allen Fällen glauben sie unschuldsvoll die Wahrheit selber auf ihrer Seite zu haben, sollte man’s für möglich halten? Und noch mehr, man glaubt es ihnen! man staunt andächtig, dass bei dieser Erziehung überhaupt etwas herauskam, wenn auch das Gegentheil des Beabsichtigten, man geht auf hundert Schritt der Vermuthung aus dem Wege, dass die unverstümmelte, „unverbesserte“ Natur doch um einige Grade reicher und schöner sein müsse als jene formalistisch zugerichtete Halbnatur. Reicher und schöner! was liegt daran? die Hauptsache ist und bleibt doch die Behaglichkeit! Dass nur der liebe Mitmensch keinen Schaden nehme! dass in den Mechanismus der allgemeinen Ordnung und Wohlfahrt nur keine von den Elementargewalten der Menschlichkeit hineinplatze! Dazu wird gestraft, unschädlich gemacht, eingesperrt, hingerichtet, dazu wird der Pflichtbegriff eingeschärft, Entsagung, Selbstverleugnung, Unterdrückung der „bösen Triebe“ gepredigt — dies ganze ausgeklügelte System von Vereinfachung, Beraubung, Castration, „Verbesserung“ des Menschen und der Natur hat keinen andern Zweck als das Glück der Völker zu sichern, die unterste Art Glück, die es giebt, als utilitarische Interessen zu fördern, die unterste Art Interessen, die es giebt, als einen bequemen behaglichen selbstgefälligen Normalmenschen in Scene zu setzen, die unterste Art Mensch, die es giebt! Heisst das nicht absichtlich Sinn und Rang aller Dinge verdrehen und Kupfer mit Golde aufwägen? heisst das nicht die That des Prometheus ungeschehen machen und dem Menschen das Feuer stehlen, damit sich die Kinder nicht wehthun?
\par

Nun, solange unsere Civilisation diese destructiven Massregeln unter Erziehung versteht, solange sie das Individuum nur auspresst, abschwächt, und das aufständische Individuum einfach opfert, beseitigt, solange sie unbequeme Energieformen vernichtet statt sie zu verwandeln, solange es weder in ihrem Willen noch Vermögen liegt, positive Eigenschaften anzuzüchten und den menschlichen Typus auch nur um eine einzige neue Function zu bereichern, — solange muss der Einzelne sich wehren, wie er kann, und darum, meine Kinder! gelobe ich, euch ein streng unmoralischer Vater zu sein! Wir wollen nicht mit offenen Augen unsere Verarmung anstreben und unserer viertausendjährigen Hochcultur ablernen, wie man auf Deficit wirthschaftet! Wir werden keinen Theil von uns durchstreichen, verneinen, verkümmern lassen; dazu wissen wir viel zu wenig von der inneren Oekonomie des Lebens, um irgend etwas, was wie Realität aussieht, fortzuwerfen. Ferne bleibe uns die Tugend, die da spricht: Führe uns nicht in Versuchung! wie sollten wir als Darsteller dessen, was wir sind, der „experimentellen Methode“ entrathen können? Nein, was in uns ist, soll auch heraus, zum mindesten ans Tageslicht der Erkenntniss; wir wollen uns nicht verschweigen und unversucht lassen. Unserem Guten und unserem Schlimmen soll die gleiche Zukunft offenstehen: „auch das Böseste ist zum Würzen würdig und zum letzten Überschäumen“. Wir wollen den Sinn des Lebens besser zu errathen suchen als jene zerstörungsgrimmigen Halbwisser, die in allem, was für den Augenblick befremdet und wehethut, eine Umsturzgefahr für das gesamte Dasein wittern und mit ihrem myopischen Verstand gegen alles, was für sie Karthago ist, ein blindwüthiges \Foreign{ceterum censeo} donnern. Wir wollen die Dinge nicht „verbessern“, wenn dazu so mörderische Eingriffe ins blühende Leben der Dinge nothwendig sind, wie jene catonischen Moraltyrannen sie ausüben und anempfehlen. Auch wir werden uns nicht immer mit der unmittelbaren Natur und Natürlichkeit begnügen; aber wenn wir den Dingen zu Leibe gehen, so halten wir uns an die Technik des Umformens, Vertheilens, Anordnens, Stilisirens, nicht an die Mechanik des Beseitigens, Ausrottens, Verstümmelns. Wir wollen den Reichthum auf Erden mehren, nicht mindern, den Reichthum an Realität, an Wärme, Klang, Farbe, Kraft, Bewegung, an Willen und Leidenschaft, sei sie gut oder böse! Die Welt soll, wenn wir scheiden, weltlicher geworden sein, positiver, gedrängter, kernhafter, zum Leben verlockender als ehedem, wenigstens für uns und Unseresgleichen! Was sonst für uns und Andere bei dieser Weltbejahung, Weltbereicherung, Weltverwirklichung abfällt, kümmert das uns? Vielleicht stossen wir uns die Köpfe ein in unserer moralisch verbalkten und verbretterten Aussenwelt: kümmert das uns? Vielleicht auch gehen wir in anderer Weise zum Teufel: vielleicht werden wir selbst moralisch? Wie, wenn jenes „biogenetische Grundgesetz“, das Goethe den Darwinisten vorentdeckt hat, auch an uns Recht behielte, wenn wir im Kleinen die Entwicklung der Menschheit recapituliren und aus Individuen zu Normalmenschen werden müssten, so wie die „blonde Bestie“ nicht umhin konnte, zu guterletzt moralischer Staats-Spiessbürger zu werden? Nun, wer wollte bei der „Gebrechlichkeit der Welt“ und dem Übergewicht der Gattung über das Individuum diese Möglichkeit schlechthin leugnen? Aber damit, meine Kinder! bin ich als euer künftiger Vater und Erzieher um so mehr im Rechte, wenn ich euch mit allem moralischen Training möglichst verschone. Wird der civilisatorische Zwang, der über die ganze Urmenschheit Herr geworden ist, auch über euch den Sieg davontragen, so habe ich nicht nöthig, eure Niederlage meinerseits zu beschleunigen; wird die Heerde, die Gesellschaft, die Moralität selbst ihr Recht von euch fordern und zu erlangen wissen, so hüte ich mich, euch zu diesem Molochsopfer noch zu schmücken und vorzubereiten — ich, der lieber den Hammer zu fassen bekäme, um das dampfende glühende Götzenbild in Trümmer zu schlagen!
\par

